\section{三项任务的逻辑链、数据闭环与结论边界}

本文将三个问题组织为一条可复核的数据链：题目要求与数据定义，特征提取与时间统一，完整模态基线，训练集局部缺失模拟，验证集模型与阈值选择，附件3专项推理，遮挡敏感性解释，以及附件4专项推理与证据定位。每个子问题均回答四个问题：输入数据是什么、模型如何处理、用什么指标验证、输出如何支持结论。

\subsection{数据分层与信息使用边界}

附件1的100条原始视频只用于问题一的特征提取与对齐；附件2按照\texttt{train/valid/test}组织，用于问题二和问题三的训练、模型选择与最终评估；附件3和附件4均为无标签专项测试集，只能在模型和阈值锁定后进行最终推理。标准化均值和标准差仅由训练集计算，并传递至验证集和测试集。

\begin{table}[htp!]
\centering\small\renewcommand{\arraystretch}{1.2}
\caption{本文各类数据的用途与允许输出}\label{tab:data_boundary}
\begin{tabularx}{\textwidth}{|l|c|X|X|}\hline
\textbf{数据} & \textbf{标签} & \textbf{用途} & \textbf{允许报告的结果} \\\hline
附件1 & 有 & 问题一特征提取、对齐和完整性核验 & 样本覆盖、特征维度、有效长度、对齐记录 \\\hline
附件2训练集 & 有 & 参数学习、标准化统计量、缺失增强 & 训练损失和训练过程 \\\hline
附件2验证集 & 有 & epoch、超参数和决策阈值选择 & Accuracy、Macro-F1、MAE、Pearson $r$ \\\hline
附件2测试集 & 有 & 全部方案锁定后的最终一次评估 & 最终性能指标 \\\hline
附件3 & 无 & 局部缺失专项推理 & 极性、强度、不确定度、缺失区间和门控权重 \\\hline
附件4 & 无 & 完整模态解释专项推理 & 极性、强度、遮挡敏感性、关键时间和证据片段 \\\hline
\end{tabularx}
\end{table}

\subsection{问题一：从原始视频到统一时序特征}

问题一的输入是视频、文本和标签三者的一一对应关系。文本经BERT编码得到$768$维上下文表示；音频由MFCC、Chroma、Mel、频谱对比度和Tonnetz等声学统计量组成$74$维特征；视觉流由逐帧视觉统计量组成$35$维特征。音频和视觉的原始有效长度保存在汇总文件中，便于核验截断与填充。

设文本有效长度为$L_t$，音频和视觉原始序列为$X_a,X_v$，统一最大长度为$50$，目标对齐长度为$L=\min(L_t,50)$。对模态$m$映射到文本时间轴的区间$I_k$，采用分段平均完成软对齐：
\begin{equation}
\widetilde X_m[k]=\frac{1}{|I_k|}\sum_{i\in I_k}X_m[i],\qquad m\in\{a,v\}.
\end{equation}
空区间使用相邻有效位置补偿，最后补零至50步，过长部分截断。问题一的交付重点不是预测准确率，而是100条样本的可追溯性：样本编号、原始路径、文本有效长度、音频/视觉有效长度、对齐后形状、标签和异常处理记录必须能够从\texttt{problem1\_feature\_summary\_100.csv}核验。

\subsection{问题二：局部缺失、鲁棒预测与专项推理}

训练阶段从完整训练样本构造连续缺失区间，缺失仅改变输入特征和对应mask，不改变标签。模型首先将三种模态投影到统一隐空间，再使用Bi-GRU提取模态内上下文，通过缺失有效mask约束跨模态注意力，最后用门控权重融合三种模态：
\begin{equation}
z=g_tz_t+g_az_a+g_vz_v,\qquad g_t+g_a+g_v=1.
\end{equation}
分类和回归输出分别为
\begin{equation}
\widehat y_c=\arg\max_j\operatorname{softmax}(W_cz)_j,\qquad \widehat y_r=f_r(z).
\end{equation}
验证集用于选择最佳epoch、缺失实验配置和BUDT阈值；测试集不参与训练过程中的方案选择。完整输入与缺失输入的性能差异定义为
\begin{equation}
\Delta M=M_{\mathrm{complete}}-M_{\mathrm{missing}}.
\end{equation}
消融结果按缺失模态、缺失位置和缺失时长分别报告Accuracy、Macro-F1、MAE和Pearson $r$，并同时给出缺失比例。附件3无标签，只汇总预测极性、预测强度、缺失区间、不确定度和三模态门控权重，不计算准确率或F1。

\subsection{问题三：遮挡敏感性与局部证据定位}

问题三复用已锁定的预测模型，对完整三模态输入进行可复核扰动。设完整输入回归输出为$\widehat y$，遮挡模态$m$后的输出为$\widehat y^{(-m)}$，定义模态遮挡敏感性为
\begin{equation}
\Delta_m=|\widehat y-\widehat y^{(-m)}|,\qquad C_m=\frac{\Delta_m}{\Delta_t+\Delta_a+\Delta_v+\varepsilon}.
\end{equation}
其中$C_m$表示相对作用程度，而不是严格独立的因果效应，因为三种模态存在交互。在作用程度最大的模态内逐时间步遮挡，定义局部敏感性$S(t)=|\widehat y-\widehat y^{(-t)}|$，再使用连续窗口求和选择最大证据窗口，并映射到原始视频秒数、帧区间和文本片段。

附件4每条记录至少保存预测极性、强度、三模态遮挡敏感性、作用比例、主要模态、显著性序列、关键时间步、视频帧和文本证据。论文中的解释应称为遮挡敏感性或反事实扰动贡献，不能直接表述为严格因果效应。

\subsection{数据闭环与结论边界}

论文中的每一个数字都应能追溯到对应的CSV、JSON或PNG输出。附件3和附件4的专项结果只能作为无标签推理展示，不能作为性能验证；门控权重只能解释为模型内部的动态分配，不等价于真实信息量。最终结论应同时报告模型性能、缺失条件、解释方法及其局限性，避免以单一最优指标代替完整实验链。
